% Target: rmp-iv-intersection-rhs-one
% Formula: sum B^{i_2} C^{i_3} R(i_1) |i_1 i_2 i_3>
% Ink: One uncovered wire followed by B--C closes through boundary tensor R.
% Boundary: Both virtual ends close into R; B, C, and R retain physical legs.
% Source: paper TN-Review-main.tex line 2054; author source ImagesReview Section
%   IV/ImagesReview Section IV.tex lines 71-78 (Rtensor)
% Capabilities: free-graph, closure, typed-ports
\[
  \begin{tenkzfree}
    \tnput[box, ports={west:virtual,east:virtual,north:physical}]
      {b}{(0mm,0mm)}{B}
    \tnput[box, ports={west:virtual,east:virtual,north:physical}]
      {c}{(8mm,0mm)}{C}
    \tnput[pill, ports={west:virtual,east:virtual,135:physical}]
      {r}{(0mm,-8mm)}{\qquad R\qquad}
    \tnput[boundary, ports={south:physical}]{rp}{(-8mm,9mm)}{}
    \tnput[boundary, ports={south:physical}]{bp}{(0mm,9mm)}{}
    \tnput[boundary, ports={south:physical}]{cp}{(8mm,9mm)}{}
    \tnjoin{b.east}{c.west}
    \tnjoin[route=arc,out=180,in=180]{b.west}{r.west}
    \tnjoin[route=arc,out=0,in=0]{c.east}{r.east}
    \tnjoin[route=hv]{r.135}{rp.south}
    \tnjoin{b.north}{bp.south} \tnjoin{c.north}{cp.south}
  \end{tenkzfree}
\]
